\documentclass[article, 12pt]{article}
\usepackage[utf8]{inputenc}
\usepackage[japanese]{babel}
\usepackage{amsmath, amssymb}
\usepackage{geometry}
\geometry{a4paper, margin=25mm}

\title{MEKOU：物理的拘束条件（ECS/メタプロトコル）による\\LLMの記号接地と実社会実装へのアプローチ}
\author{高久}
\date{2026年4月15日}

\begin{document}

\maketitle

\section{研究の背景と動機}
2040年問題に象徴される深刻な労働力不足、特に介護・医療等のブルーカラー領域における供給不足は喫緊の課題である。従来のロボティクス（例：Spot等）は、End-to-Endの強化学習による「条件反射」の獲得には長けているが、未知のタスクに対する「意味論的なプランニング（エージェント能力）」が欠如している。本研究では、LLM（大規模言語モデル）の持つ意味空間を、物理的拘束条件下での最適化問題として再定義することで、この課題を解決する。

\section{提案アーキテクチャ：MEKOU Engine}
本システムは、以下の2つの主要コンポーネントを核としたパイプラインによって構成される。

\begin{itemize}
    \item \textbf{ECS (Entity Component System):} 
    現実世界の文脈を「エンティティ」と「コンポーネント」の集合 $A$ として階層的に定義する。これにより、水平スケーリングと疎結合なシステム拡張を実現し、現実の物理状態をデータ構造として厳密に管理する。
    
    \item \textbf{メタプロトコル:} 
    自然言語（テキスト）をインターフェースとし、LLMとECSを直結させる。IoTデバイス側の物理制限（例：回転角の閾値）をハードコードされた制約として定義し、実行不可能な命令を動的に排除する。
\end{itemize}

\section{数理的仮説：拘束条件付き最適化問題としてのLLM}
LLMを単なる確率機ではなく「意味空間を持った行列」と捉える。入力プロンプトを目的関数 $f(t)$ とし、ECSおよびメタプロトコルによる物理制約を $g(x)$ としたとき、AIの推論は以下のラグランジュ未定乗数法的な最適化問題に帰着できる。

\begin{equation}
\min \quad \mathcal{J} = \| \text{Input} - \text{LLM}(t) \|^2 \quad \text{subject to} \quad g_{\text{ECS/Meta}}(x) \leq 0
\end{equation}

ここで、RAG（Retrieval-Augmented Generation）を「時間依存のタイムスタンプ」として機能させることで、現代制御工学における動的システムの最適化と同様の枠組みで、AIを現実空間へ「接地（Grounding）」させる。

\section{エンジニアリングによる記号接地問題の解決}
記号接地問題における「記号」と「現実」の乖離を、メタプロトコルによるハードコードされた拘束によって強制的に接続する。これは「サイコロの振り方（アルゴリズム）」を研究するのではなく、「サイコロの形状（構造）」そのものを設計し、おもりを仕込むことで、必然的に正解（現実的な回答）を出力させるアプローチである。

\section{結論}
本研究は、低レイヤーのエンジニアリング（Rustによるエンジン実装）から出発し、エージェントの本質を「賢さ」ではなく「物理的制約」に見出したものである。この構造化されたパイプラインにより、安全かつ決定論的なアンドロイド開発およびIoTオーケストレーションの基盤を提供する。

\end{document}